\documentclass[11pt,a4paper]{article}
\usepackage[T1]{fontenc}
\usepackage[utf8]{inputenc}
\usepackage[english]{babel}
\usepackage{lmodern,amsmath,amssymb,amsthm,booktabs,array}
\usepackage[margin=26mm]{geometry}
\usepackage{microtype}
\setlength{\emergencystretch}{2em}
\usepackage[colorlinks=true,linkcolor=blue,citecolor=blue,urlcolor=blue]{hyperref}
\hypersetup{pdftitle={A degree-two neighborhood bound for total domination in trees},
 pdfauthor={Jonas J. Gebendorfer},pdfsubject={Total domination; WOWII 348},
 pdfkeywords={total domination, trees, degree-two vertices, Graffiti.pc}}
\newtheorem{theorem}{Theorem}[section]
\newtheorem{lemma}[theorem]{Lemma}
\newtheorem{corollary}[theorem]{Corollary}
\newtheorem{proposition}[theorem]{Proposition}
\theoremstyle{remark}
\newtheorem{remark}[theorem]{Remark}
\newcommand{\td}{\gamma_t}
\newcommand{\BB}{B}
\newcommand{\ind}{\mathbf{1}}
\newcommand{\parv}{\operatorname{par}}
\newcommand{\ch}{\operatorname{Ch}}
\newcommand{\OR}{\mathbin{\vee}}
\title{A degree-two neighborhood bound\newline for total domination in trees}
\author{Jonas J. Gebendorfer}
\date{13 October 2026}
\begin{document}
\maketitle

\begin{abstract}
Let $T$ be a finite simple tree of order at least three, let $D_2(T)$ be
the set of its degree-two vertices, and put
$b(T)=|N_T(D_2(T))\setminus D_2(T)|$, where $N_T$ denotes the open
neighborhood. We prove that every total dominating set $D$ of $T$ satisfies
$3|D|\ge 2b(T)+2$. Consequently,
$\gamma_t(T)\ge\lceil(2b(T)+2)/3\rceil$, which implies the inequality
recorded as Conjecture 348 in \emph{Written on the Wall II}.
The proof is a structural induction using seven explicit bounds for
rooted subtrees, with no upper bound on the order or degrees.
The additive constant is best possible with the coefficients $3$ and $2$
fixed, as witnessed by $P_3$ and $P_4$. An explicit family of connected
graphs shows that the tree hypothesis is essential even for the original
inequality. A supplementary standard-library Python package checks all
labelled trees of orders three through seven and includes a second
algorithm for exact total domination. These finite checks are separate
from the proof of the universal statement.
\end{abstract}

\noindent\textbf{Keywords.} Total domination; trees; degree-two vertices;
Graffiti.pc; rooted-tree induction.\\
\textbf{2020 Mathematics Subject Classification.} 05C69, 05C05.

\section{Introduction and statement}
\label{sec:introduction}

Graffiti.pc generated a number of conjectures concerning the total
domination number of a graph. DeLaVi\~na's collection \emph{Written on
the Wall II} \cite{wowii} records the following assertion as number 348,
dated 18 February 2009: if $T$ is a tree on $n>2$ vertices, then
\begin{equation}\label{eq:original}
 \td(T)\ge\frac{2}{3}\,|N_T(D_2(T))\setminus D_2(T)|,
 \qquad D_2(T)=\{v\in V(T):\deg_T(v)=2\}.
\end{equation}
The original entry uses $\gamma_T$ for the total domination number;
we use the customary notation $\gamma_t$ throughout. The neighborhood
and set subtraction in \eqref{eq:original} are those of the original
definitions. In particular, degree-two vertices lying in
$N_T(D_2(T))$ do not contribute to the quantity on the right.

The study of Graffiti.pc bounds for total domination in trees includes
the paper of DeLaVi\~na, Larson, Pepper, and Waller \cite{dlpw}.
Desormeaux and Henning \cite{dh} also investigated Graffiti.pc lower
bounds for total domination. We supply here a direct proof of
\eqref{eq:original} with an additional constant. No result from either
paper is used as a premise. The source identification and the limits of
the literature comparison are stated in Appendix~\ref{app:sources};
no assertion of first priority is made.

All graphs in this paper are finite, simple, and undirected. A tree is
a connected graph containing no cycle. Its order is $|V(T)|$.
For a graph $G$ and a vertex $v$, let
\[
 N_G(v)=\{w\in V(G):vw\in E(G)\},\qquad \deg_G(v)=|N_G(v)|.
\]
For $X\subseteq V(G)$, define its open neighborhood by
\[
 N_G(X)=\bigcup_{x\in X}N_G(x).
\]
This union can intersect $X$; a vertex does not belong to the union
merely because it belongs to $X$. A set $D\subseteq V(G)$ is
\emph{total dominating} if $N_G(v)\cap D\ne\varnothing$ for every
$v\in V(G)$. Thus every vertex of $D$ itself must have a neighbor in
$D$. The total domination number $\td(G)$ is the minimum size of such
a set. It exists for any graph with no isolated vertex, since its full
vertex set is total dominating. A tree of order at least two has no
isolated vertex.

For a tree $T$, set
\begin{equation}\label{eq:boundary}
 D_2(T)=\{v\in V(T):\deg_T(v)=2\},\qquad
 \BB(T)=N_T(D_2(T))\setminus D_2(T),\qquad b(T)=|\BB(T)|.
\end{equation}
Equivalently, $v\in\BB(T)$ if and only if $\deg_T(v)\ne 2$ and
$v$ has a neighbor of degree two. This follows directly
from the open neighborhood and the set difference.

\begin{theorem}\label{thm:main}
Let $T$ be a tree of order at least three. Every total dominating set
$D$ of $T$ satisfies
\begin{equation}\label{eq:main}
 3|D|\ge 2b(T)+2.
\end{equation}
\end{theorem}

\begin{corollary}\label{cor:348}
Every tree $T$ of order at least three satisfies
\begin{equation}\label{eq:ceiling}
 \td(T)\ge\left\lceil\frac{2b(T)+2}{3}\right\rceil.
\end{equation}
In particular, \eqref{eq:original} holds on its full stated domain.
\end{corollary}

\begin{proof}
Apply Theorem~\ref{thm:main} to a minimum total dominating set.
Division by three gives $\td(T)\ge(2b(T)+2)/3$. Since
$\td(T)$ is an integer, this implies \eqref{eq:ceiling}. Finally,
$(2b(T)+2)/3\ge 2b(T)/3$, proving \eqref{eq:original}.
\end{proof}

\section{The rooted-subtree bounds}
\label{sec:local}

Fix a tree $T$ of order at least three, a total dominating set $D$,
and initially an arbitrary root $r$. Between any two vertices of a
tree there is a unique simple path: connectedness gives a path, and
two different simple paths would contain a cycle. For $v\ne r$, its
\emph{parent} $\parv(v)$ is the neighbor on its path to $r$.
The other neighbors are its \emph{children}, whose set is
$\ch(v)$. A nonroot vertex of degree $d$ has exactly $d-1$
children. Let $T_v$ be the subtree induced by $v$ and its descendants,
including $v$. The sets $V(T_w)$ for $w\in\ch(v)$ are pairwise
disjoint and partition $V(T_v)\setminus\{v\}$, because each descendant
path to $v$ passes through a unique first child.

For every vertex $x$, put $s_x=\ind_{\{x\in D\}}$. For $v\ne r$, put
\[
 p_v=s_{\parv(v)},\qquad
 q_v=\ind_{\{\deg_T(\parv(v))=2\}}.
\]
Define the subtree weight
\begin{equation}\label{eq:weight}
 F(v)=\sum_{x\in V(T_v)}\bigl(3s_x-2\ind_{\{x\in\BB(T)\}}\bigr).
\end{equation}
All degrees, and all membership tests for $\BB(T)$, refer to the
original tree $T$. They are never recomputed in $T_v$.
The parent of $v$ is outside the sum in \eqref{eq:weight}.
The disjoint partition above gives the exact identity
\begin{equation}\label{eq:recurrence}
 F(v)=3s_v-2\ind_{\{v\in\BB(T)\}}+
       \sum_{w\in\ch(v)}F(w).
\end{equation}

Two consequences of total domination will be used repeatedly. If
$v\ne r$ and $p_v=0$, then $v$ has at least one selected child:
its neighbors are its parent and children, and the parent is not in
$D$. This holds whether or not $s_v=1$. A nonroot leaf must
have $p_v=1$, since its parent is its only neighbor.

\begin{lemma}\label{lem:states}
For every nonroot vertex $v$, the applicable row of Table~\ref{tab:states}
is a lower bound for $F(v)$. A nonroot leaf with $p_v=0$ is impossible.
\end{lemma}

\begin{table}[htbp]
\centering
\renewcommand{\arraystretch}{1.15}
\begin{tabular}{@{}llll@{}}
\toprule
$\deg_T(v)$ & $s_v$ & Parent condition & Bound for $F(v)$\\
\midrule
$1$ & $0$ or $1$ & $p_v=1$ & $3s_v-2q_v$ (exact)\\
$2$ & $0$ & Any & $2$\\
$2$ & $1$ & $p_v=0$ & $4$\\
$2$ & $1$ & $p_v=1$ & $1$\\
$\ge 3$ & $0$ & Any & $2$\\
$\ge 3$ & $1$ & $p_v=q_v=1$ & $1$\\
$\ge 3$ & $1$ & All other cases & $2$\\
\bottomrule
\end{tabular}
\caption{Lower bounds in the original tree. \emph{Any} allows either value
of each unspecified parent indicator, subject to total domination.}
\label{tab:states}
\end{table}

\begin{proof}
We induct on the height of $T_v$, the maximum number of edges
on a path from $v$ to a descendant. A nonroot vertex of
height zero has no child and therefore degree one. At positive
height, every child has smaller height, so all applicable table rows
are already valid for that child. Each case below uses only
these previously established child bounds and \eqref{eq:recurrence}.

\smallskip\noindent\emph{Case 1: $\deg_T(v)=1$.}
The only neighbor of $v$ is its parent. Hence $p_v=1$.
Since $v\notin D_2(T)$, it belongs to $\BB(T)$ precisely when its
parent has degree two, that is, when $q_v=1$. It has no child, and
\eqref{eq:weight} gives $F(v)=3s_v-2q_v$. This proves
the induction base and the first row.

\smallskip\noindent\emph{Case 2: $\deg_T(v)=2$.}
There is exactly one child $w$. The subtraction in
\eqref{eq:boundary} ensures $v\notin\BB(T)$. Thus
\begin{equation}\label{eq:degree2rec}
 F(v)=3s_v+F(w),\qquad (p_w,q_w)=(s_v,1).
\end{equation}
There are three selection conditions at $v$.

If $s_v=0$, then $p_w=0$. The child $w$ cannot be a leaf: its only
neighbor would be the unselected vertex $v$. For a degree-two child
with $p_w=0$, the induction hypothesis gives $F(w)\ge 2$ if
$s_w=0$, and $F(w)\ge 4$ if $s_w=1$. For a child of degree at least
three, it gives $F(w)\ge 2$ when $s_w=0$, and also $F(w)\ge 2$
when $s_w=1$, since $p_w=0$ excludes the sixth row.
These are all possible child degrees. Therefore $F(w)\ge 2$, and
\eqref{eq:degree2rec} gives $F(v)\ge 2$.

If $s_v=1$ and $p_v=0$, total domination of $v$ forces $s_w=1$.
Now $(p_w,q_w)=(1,1)$. A selected leaf has weight $3-2=1$ by
Case~1. A selected degree-two child has weight at least $1$ by the
fourth row of the induction hypothesis. A selected child of degree
at least three has weight at least $1$ by its sixth row. Hence
$F(w)\ge 1$ and $F(v)=3+F(w)\ge 4$.

If $s_v=1$ and $p_v=1$, again $(p_w,q_w)=(1,1)$, but $s_w$ may
have either value. The possible child bounds are $-2$ for an
unselected leaf, $1$ for a selected leaf, $2$ for an unselected
degree-two child, $1$ for a selected degree-two child, $2$ for an
unselected child of degree at least three, and $1$ for a selected
child of degree at least three. All are at least $-2$.
Consequently $F(v)=3+F(w)\ge 1$. This proves the three degree-two rows.

\smallskip\noindent\emph{Case 3: $\deg_T(v)\ge 3$ and $s_v=0$.}
The vertex $v$ has at least two children. Each child $w$ has
$(p_w,q_w)=(0,0)$, since $v$ is unselected and not of degree two.
No child can be a leaf, by total domination of that child.
For a degree-two child, the second and third rows give bounds $2$
and $4$, according to its selection. For a child of degree at least
three, the fifth and seventh rows both give $2$; the sixth row
cannot apply because $p_w=0$. Thus every child contributes at least
$2$. The contribution of $v$ itself is at least $-2$, and so
\[
 F(v)\ge -2+2|\ch(v)|\ge -2+2\cdot 2=2.
\]
There is no upper bound on the number of children.

\smallskip\noindent\emph{Case 4: $\deg_T(v)\ge 3$ and $s_v=1$.}
Each child $w$ has $(p_w,q_w)=(1,0)$. The induction hypothesis
gives this exhaustive list of child contributions:
\[
\begin{array}{c|cc}
 &s_w=0&s_w=1\\ \hline
\deg_T(w)=1&0&3\\
\deg_T(w)=2&\ge 2&\ge 1\\
\deg_T(w)\ge 3&\ge 2&\ge 2
\end{array}
\]
Every child weight is nonnegative, every selected child has weight
at least $1$, and every degree-two child has weight at least $1$.

If $q_v=1$, the parent of $v$ has degree two. Since $v$ does
not have degree two, $v\in\BB(T)$, and its contribution is
$3-2=1$. If $p_v=1$, the nonnegative child sum gives $F(v)\ge 1$,
as in the sixth row. If $p_v=0$, total domination of $v$ requires
a selected child. That child contributes at least $1$, so
$F(v)\ge 1+1=2$, as in the seventh row.

If $q_v=0$, the parent of $v$ does not have degree two.
Therefore $v\in\BB(T)$ if and only if at least one child has degree
two, since its parent and children are all its neighbors.
When such a child exists, $v$ contributes $1$, that child contributes
at least $1$, and all other contributions are nonnegative.
Hence $F(v)\ge 2$. Otherwise $v\notin\BB(T)$, its contribution is
$3$, and $F(v)\ge 3\ge 2$. This proves the seventh row for either
value of $p_v$.

The cases cover every nonroot vertex and every feasible selection
of its indicators. The induction is complete.
\end{proof}

\section{Closing the induction at a leaf}
\label{sec:root}

\begin{proof}[Proof of Theorem~\ref{thm:main}]
A finite tree of order at least two has a leaf. Choose a longest
simple path. An endpoint has no neighbor outside the path, since
that would extend it, and no neighbor on the path other than the
next vertex, since that would create a cycle. Its degree is one.

Choose a leaf $r$ as the root and let $v$ be its unique neighbor.
Total domination of $r$ forces $s_v=1$. Since $|V(T)|\ge 3$,
we have $\deg_T(v)\ge 2$: if both endpoints of $rv$ had degree one,
connectedness would force $V(T)=\{r,v\}$.
Also $p_v=s_r$ and $q_v=0$, since $\deg_T(r)=1$.
All vertices except $r$ belong to $T_v$, and therefore
\begin{equation}\label{eq:globalweight}
 3|D|-2b(T)=3s_r-2\ind_{\{r\in\BB(T)\}}+F(v).
\end{equation}

If $\deg_T(v)=2$, then $r\in\BB(T)$, since $r$ is a leaf adjacent
to a degree-two vertex. When $s_r=0$, the third row of
Table~\ref{tab:states} gives $F(v)\ge 4$, and
\[
 3|D|-2b(T)\ge -2+4=2.
\]
When $s_r=1$, the fourth row gives $F(v)\ge 1$, and
\[
 3|D|-2b(T)\ge (3-2)+1=2.
\]

If $\deg_T(v)\ge 3$, then $r\notin\BB(T)$, because its only
neighbor does not have degree two. Since $s_v=1$ and $q_v=0$,
the seventh row gives $F(v)\ge 2$ for either value of $p_v=s_r$.
Equation~\eqref{eq:globalweight} yields
\[
 3|D|-2b(T)\ge 3s_r+2\ge 2.
\]
These are all possibilities for $v$. The argument applies to the
arbitrary total dominating set $D$, proving \eqref{eq:main}.
\end{proof}

\begin{remark}
The proof does not require $D$ to be minimum, inclusion-minimal, or
free of leaves. The seven local bounds apply to the actual selection
of any total dominating set; no exchange argument is needed.
\end{remark}

\section{Equality examples and the tree hypothesis}
\label{sec:examples}

Let $P_m$ be the path with vertices $z_1,\ldots,z_m$ and edges
$z_i z_{i+1}$ for $1\le i<m$.

\begin{proposition}\label{prop:constant}
The additive constant $2$ in \eqref{eq:main} is best possible among
uniform bounds with the coefficients $3$ and $2$ fixed. Equality
holds for $P_3$ and $P_4$ with a minimum total dominating set.
\end{proposition}

\begin{proof}
For $P_3$, the only degree-two vertex is $z_2$, and
$\BB(P_3)=\{z_1,z_3\}$. The set $\{z_1,z_2\}$ is total dominating:
$z_1$ and $z_2$ are adjacent to one another, and $z_3$ is adjacent
to $z_2$. No singleton is total dominating in a simple graph,
because its selected vertex has no selected neighbor. Hence
$\td(P_3)=2$ and $3\td(P_3)-2b(P_3)=6-4=2$.

For $P_4$, $D_2(P_4)=\{z_2,z_3\}$ and $N(D_2(P_4))=V(P_4)$.
Subtracting $D_2(P_4)$ gives $\BB(P_4)=\{z_1,z_4\}$.
The set $\{z_2,z_3\}$ is total dominating; the singleton obstruction
gives $\td(P_4)=2$. Again $3\td(P_4)-2b(P_4)=2$.
Either example excludes a uniform additive constant greater than $2$.
\end{proof}

\begin{remark}
These examples establish sharpness of the additive constant with
the displayed coefficients fixed. They do not establish asymptotic
optimality of $2/3$ as $b(T)$ grows. Trees without degree-two vertices
have $b(T)=0$ and are included in the theorem. The tree $K_2$ is
outside the domain of \eqref{eq:original}, but satisfies
\eqref{eq:main} as well: $b(K_2)=0$ and $\td(K_2)=2$.
The one-vertex tree has no total dominating set.
\end{remark}

In the following proposition, $D_2(G)$ and $\BB(G)$ have the
definitions in \eqref{eq:boundary} with $G$ replacing $T$.

\begin{proposition}\label{prop:cycles}
For every integer $k\ge 3$, there is a connected simple graph $G_k$
of order $3k+2$ with $\td(G_k)=2$ and $|\BB(G_k)|=k+1$.
Even \eqref{eq:original} consequently fails if its tree hypothesis
is replaced by connectedness.
\end{proposition}

\begin{proof}
Take distinct vertices
\[
 V(G_k)=\{u,v\}\cup\{c_i,x_i,y_i:1\le i\le k\}
\]
and precisely the edges
\[
 E(G_k)=\{uv\}\cup
 \bigcup_{i=1}^{k}\{vc_i,ux_i,uy_i,c_ix_i,c_iy_i\}.
\]
All endpoints are distinct within each edge, and no edge is repeated.
Each $c_i$ is joined to $v$, each $x_i,y_i$ to $u$, and $u$ to $v$,
so the graph is simple and connected. The distinct vertices
$u,v,c_i,x_i$ form a cycle of length four.

The degrees are
\[
 \deg(u)=2k+1,\quad \deg(v)=k+1,\quad
 \deg(c_i)=3,\quad \deg(x_i)=\deg(y_i)=2.
\]
For $k\ge 3$, the degree-two vertices are exactly the $x_i,y_i$.
Their neighborhoods are $\{u,c_i\}$; thus
$\BB(G_k)=\{u,c_1,\ldots,c_k\}$ and $|\BB(G_k)|=k+1$.
The set $\{u,v\}$ is total dominating: the selected vertices
dominate each other, each $c_i$ has selected neighbor $v$, and
each $x_i,y_i$ has selected neighbor $u$. The singleton obstruction
shows $\td(G_k)=2$. The extended version of \eqref{eq:original}
would require $6\ge 2(k+1)$, which is false for $k\ge 3$.
\end{proof}

\section{Supplementary exact verification}
\label{sec:verification}

The proof of Theorem~\ref{thm:main} has no computational premise.
The accompanying package supplies finite checks of the definitions,
local bounds, examples, and minimum values. We specify the checks
so that they can be reproduced without an internal workspace or
an external service.

\subsection{Representation and exhaustive checks}

Vertices are labelled $0,\ldots,n-1$. A subset is an integer mask
whose bit $i$ is one exactly when vertex $i$ is selected.
The neighborhood mask $A_i$ has a one at each neighbor of $i$.
A mask $M$ is accepted as total dominating precisely when
\begin{equation}\label{eq:masktd}
 A_i\mathbin{\&}M\ne 0\quad\hbox{for every }i,
\end{equation}
where $\&$ is bitwise intersection. This is the definition of total
domination, including at selected vertices. Degrees are the numbers
of one-bits in $A_i$. If $Q$ is the degree-two mask and $U$ the
bitwise union of its vertices' neighborhood masks, then the boundary
mask is $U\mathbin{\&}\neg Q$. Since $U$ has no bit outside
$0,\ldots,n-1$, Python's unbounded integer complement adds no vertex.

The routine \texttt{adjacency} rejects order less than three,
invalid endpoint labels, loops, repeated edges, and disconnected
inputs. In tree mode it requires exactly $n-1$ edges.
This recognizes the tree domain: from one vertex in a connected
graph, repeatedly choose an edge to a new vertex. These edges form
a connected acyclic spanning subgraph with $n-1$ edges. If the
original graph has $n-1$ edges, it equals that subgraph.
Conversely, leaf deletion and induction show that every tree of
order $n$ has $n-1$ edges.

For each $n=3,\ldots,7$, the program decodes every length-$(n-2)$
sequence over $\{0,\ldots,n-1\}$ by the Pr\"ufer procedure.
Appendix~\ref{app:pruefer} proves that this gives every labelled tree
exactly once. For each tree, all $2^n$ masks are tested. Every accepted
total dominating set $D$ is checked for $3|D|\ge 2b(T)+2$.
The minimum accepted size is compared with the four-state optimizer
of Appendix~\ref{app:dp}. At orders three through five, every accepted
set is also checked against every nonroot bound in
Table~\ref{tab:states}, under every possible root, using original-tree
degrees and boundary membership.

\begin{table}[htbp]
\centering
\small
\setlength{\tabcolsep}{6pt}
\begin{tabular}{@{}rrrrr@{}}
\toprule
Order & Labelled trees & Subsets tested & Accepted TD sets & Equality trees\\
\midrule
3 & 3 & 24 & 9 & 3\\
4 & 16 & 256 & 76 & 12\\
5 & 125 & 4\,000 & 855 & 60\\
6 & 1\,296 & 82\,944 & 13\,626 & 120\\
7 & 16\,807 & 2\,151\,296 & 277\,431 & 210\\
\midrule
Total & 18\,247 & 2\,238\,520 & 291\,997 & 405\\
\bottomrule
\end{tabular}
\caption{Exhaustive labelled-tree checks. An equality tree satisfies
$3\td(T)=2b(T)+2$; different labellings are counted separately.}
\label{tab:counts}
\end{table}

Table~\ref{tab:counts} gives the frozen expected output.
There are $18\,066$ further local-bound checks under all roots at
orders three through five. A further $70$ instances consist of a
path, a star, and twelve Pr\"ufer trees for each order in
$\{8,12,20,50,100\}$. One generator \texttt{random.Random(20261013)}
is used, visiting these orders in the displayed order; each code is
obtained by $n-2$ calls to \texttt{randrange(n)}. These larger cases
use exact optimization and check the selected-set witnesses.
They constitute a finite sample.

\subsection{Calibration and reproducibility}

Before the enumeration, calibration tests check the minimum values
$2,2,3,4,4$ for $P_3$ through $P_7$ by subset enumeration and the
optimizer. They accept an adjacent pair on $P_3$, reject its middle
singleton, and reject $\{z_1,z_3\}$ on $P_4$. They check the
degree-two and boundary masks of $P_3$ and $P_4$. On $P_6$,
$|N(D_2)|=6$ but $|\BB|=2$: omitting the subtraction would
incorrectly demand $3\cdot 4\ge 2\cdot 6+2$.

The local checker is deliberately given a damaged bound function
that replaces the unselected-leaf value $-2$ by $-1$ when the parent
has degree two. For the actual $P_3$ selection $\{z_1,z_2\}$, root
$z_1$ makes $z_3$ an unselected leaf of weight $-2$; the checker
must reject the damaged bound. Invalid graph inputs exercise the
domain rejections. The eleven-vertex $G_3$ is accepted as a connected
graph and rejected as a tree. For orders three through five, the
Pr\"ufer edge sets are compared with the connected graphs obtained
by enumerating every $(n-1)$-element subset of all possible edges;
the two sets coincide.

The archive includes the complete LaTeX source, its compiled PDF,
the verifier and replay driver, frozen expected results, source
metadata, and a file-digest manifest. The computational files are
\begin{quote}\small\ttfamily
scripts/verify348.py\quad reproduce.py\\
evidence/verification348.json\quad manifest348.json
\end{quote}
Python 3.10 or later, with optimization
disabled, suffices for verification; no third-party Python package
or network is used. From the extracted directory, run
\begin{quote}\ttfamily
python -B -X utf8 reproduce.py
\end{quote}
The driver verifies every file digest listed in the SHA-256 manifest,
copies those files to a fresh temporary directory, runs the verifier
there, and compares the entire output JSON with the frozen report.
A mismatch or failed check gives a nonzero exit status. Success
prints \texttt{PASS} and the aggregate counts. Assertions must remain
enabled; the verifier explicitly refuses optimized execution.
The README gives direct verifier and LaTeX build commands.

Subset enumeration, optimization, and direct edge-set enumeration
differ in mechanism, but share graph and domination definitions in
this implementation. They are not an external independent review
or a formal proof-assistant verification. Hashes establish byte
identity relative to the supplied manifest, rather than mathematical
correctness or priority. The universal warrant is the induction of
Lemma~\ref{lem:states} and the proof of Theorem~\ref{thm:main}.

\appendix
\section{Coverage of the labelled-tree enumeration}
\label{app:pruefer}

Given a labelled tree of order $n\ge 2$, repeatedly delete its
least-labelled current leaf and record that leaf's neighbor,
stopping when two vertices remain. Leaf deletion preserves
connectedness and creates no cycle, so this is defined throughout.
It records a sequence of length $n-2$, the Pr\"ufer code; the
two surviving vertices are adjacent.

To decode a sequence $a_1,\ldots,a_{n-2}$, initialize
\[
 d_i=1+|\{j:a_j=i\}|\quad(0\le i<n).
\]
At step $j$, choose the least label $u$ with current $d_u=1$, add
$ua_j$, and decrease $d_u$ and $d_{a_j}$ by one.
Finally join the two labels with current degree value one.
Removed vertices have degree value zero.

When $m$ labels are active, the remaining code has length $m-2$.
Each active degree value is one plus the occurrences of its label
in that code, and removed labels have no remaining occurrence.
This invariant holds initially and is preserved at each step.
At least two active labels have no occurrence, because $m-2$
symbols cannot contain all but one of $m$ labels. The next symbol
$a_j$ has a remaining occurrence, so it is active with value at
least two and differs from the chosen $u$. Removing $u$ loses no
later symbol. Thus the decoder is defined at every step and ends
with exactly two active labels of value one.

Reverse the construction from its final edge. Each reverse step
attaches a new vertex by one edge to an already present vertex,
preserving connectedness and acyclicity. The result is a tree.
In the active graph at any stage, each label has one incident edge
at its own later removal or at the final pair, and one additional
incident edge per occurrence in the remaining code.
Its actual degree equals the maintained degree value. Encoding
the decoded tree therefore chooses the same least leaf at each
step and returns the starting code.

Conversely, encoding a tree records label $i$ once for each incident
edge deleted before $i$ is removed, or before the final edge remains.
Exactly one incident edge survives to that point. Hence its code
contains $i$ exactly $\deg_T(i)-1$ times.
Decoder initialization recovers the actual degrees; the first
least leaf and its recorded neighbor agree. After deleting that
leaf, the same argument applies to the smaller tree. Induction
recovers all edges, including the final edge. The procedures are
mutual inverses, so the $n^{n-2}$ sequences give exactly that many
labelled trees, with no repetition or omission.

\section{The exact four-state optimizer}
\label{app:dp}

Root a tree at any vertex $r$. For $a,b\in\{0,1\}$, let $C_v(a,b)$
be the minimum size of $S\subseteq V(T_v)$ with these conditions:
\begin{enumerate}
\item $v\in S$ if and only if $a=1$;
\item every vertex of $T_v$ other than $v$ has a neighbor in $S$;
\item $v$ has a selected child if and only if $b=1$.
\end{enumerate}
Domination of $v$ is not required at this stage; its parent may
eventually supply a neighbor. An infeasible state has cost
$+\infty$ and is omitted in the implementation.
At a leaf the values are $C_v(0,0)=0$, $C_v(1,0)=1$, and
$C_v(0,1)=C_v(1,1)=+\infty$.

Let $w_1,\ldots,w_m$ be the children of $v$ and fix $a$.
Choose feasible child states $(a_i,b_i)$. The child $w_i$ is
dominated after attaching it to $v$ exactly when
\begin{equation}\label{eq:dpcompat}
 b_i\OR a=1.
\end{equation}
Its neighbors are its children and $v$; $b_i$ records a selected
child and $a$ the selection of $v$. Its own selection is irrelevant
to its domination. Therefore
\begin{equation}\label{eq:dprec}
 C_v(a,b)=a+
 \min\left\{
 \sum_{i=1}^{m}C_{w_i}(a_i,b_i):
 b_i\OR a=1\ \hbox{for every }i,\ 
 \bigvee_{i=1}^{m}a_i=b
 \right\}.
\end{equation}
An empty disjunction is zero; a minimum with no feasible choice
is $+\infty$.

For the lower-bound direction, restrict any feasible set at $v$
to each child subtree. Every proper descendant of $w_i$ has all
its neighbors in that subtree, by uniqueness of paths, and is
therefore dominated within the restriction. It defines a feasible
child state. Domination of $w_i$ gives \eqref{eq:dpcompat}, and
the selected-child condition at $v$ gives the disjunction constraint.
Disjointness of child subtrees gives total size $a$ plus the child
sizes. Replacing each size by its state minimum proves the lower bound.

For the upper-bound direction, choose minimum sets for a feasible
collection of child states and include $v$ exactly when $a=1$.
Proper descendants are dominated by their state conditions,
children by \eqref{eq:dpcompat}, and the disjunction gives the
selected-child condition at $v$. The union is feasible with the
displayed cost. This proves \eqref{eq:dprec} in both directions.

At the root no parent can dominate it, so a selected child is necessary:
\begin{equation}\label{eq:dproot}
 \td(T)=\min\{C_r(0,1),C_r(1,1)\}.
\end{equation}
Every total dominating set restricts to one of these root states,
and every feasible root state gives a total dominating set, proving
\eqref{eq:dproot}.

The implementation processes children one at a time.
For fixed $a$, the initial cost is $a$ and the initial disjunction is
zero. A compatible child state adds its cost and changes the
running disjunction to its disjunction with $a_i$.
The least cost is retained for each resulting bit. This implements
repeated minimization of exactly the same finite expression.
Each cost carries the union of its selected-vertex masks; the final
mask is checked against \eqref{eq:masktd} and its cardinality against
the cost. All calculations are integer calculations without a
numerical optimization tolerance.

\section{Source identification and literature scope}
\label{app:sources}

The primary source identifying the target is the 23 July 2026
archival capture of DeLaVi\~na's WOWII page \cite{wowii}, retrieved
on 13 October 2026. Entry 348 is dated 18 February 2009 and has
historical marker \texttt{O}. Its original linked definitions are
94 (total domination), 46 (neighborhood of a vertex subset), and
66 (set difference); the degree-two set is defined in the entry.
These identify \eqref{eq:original} with its quantifiers and
open-neighborhood convention. The source record retains the row,
definition excerpts, URLs, retrieval dates, and SHA-256 digests.
Two later public transcriptions agree, but both derive from the
collection and are not independent mathematical reviews.

The live original curator endpoint was not reached. The archival
marker and a current secondary catalogue do not establish the
current original curator status. The full text of
Desormeaux--Henning \cite{dh} was not available in the literature
check; only its bibliographic record and abstract were inspected.
Its possible relationship to this exact inequality is unchecked.
The paper proves the recorded mathematical statement and the
displayed strengthening; worldwide first priority is not established
by the bounded source search.

\section*{Acknowledgment and assistance disclosure}
OpenAI Codex assisted in developing the proof, drafting the exposition,
and preparing the supplementary programs. It is not listed as an author.
The package records self-checks and a reproducible finite replay,
rather than external peer review.

\section*{Data and code availability}
The supplementary archive contains the complete LaTeX source,
compiled manuscript, standard-library Python verifier and replay driver,
frozen expected results, source metadata, and a SHA-256 manifest.
These files suffice to read the proof and run the verification
without an internal archive. No public repository or DOI for this
manuscript is asserted.

\begin{thebibliography}{9}
\bibitem{wowii}
E. DeLaVi\~na,
\emph{Written on the Wall II}, Graffiti.pc conjecture collection,
entry 348, dated 18 February 2009.
Original page:
\url{https://cms.uhd.edu/faculty/delavinae/research/wowII/open.html}.
\href{https://web.archive.org/web/20260723211215id_/http://cms.uhd.edu/faculty/delavinae/research/wowII/open.html}{Archival capture of 23 July 2026}.
Accessed 13 October 2026.

\bibitem{dlpw}
E. DeLaVi\~na, C. E. Larson, R. Pepper, and B. Waller,
Graffiti.pc on the total domination number of a tree,
\emph{Congressus Numerantium} \textbf{195} (2009), 5--18.
Author-hosted manuscript: \url{https://math1um.github.io/Research/r05.pdf}.

\bibitem{dh}
W. J. Desormeaux and M. A. Henning,
Lower bounds on the total domination number of a graph,
\emph{Journal of Combinatorial Optimization} \textbf{31} (2016), 52--66.
\url{https://doi.org/10.1007/s10878-014-9708-2}.
Bibliographic record and abstract inspected; full text not inspected.
\end{thebibliography}

\end{document}
